\section{Additional New-data Results}
\label{app:new-results-details}

\textbf{Aggregation and interval width.}
All public-data summaries average the saved fold or repetition metrics without weighting by fold size.
For Parkinsons, each fold metric averages recordings, so patients with more recordings contribute more within their fold.
Table~\ref{tab:new-widths} supplements the main results with mean prediction-interval widths.
Repeated row holdouts share observations, so their spread describes variation across partitions rather than independent-sample uncertainty.
The two Airfoil settings and the Superconductivity pilot and full-data runs are not paired comparisons.
The 5{,}000-row pilot is not pooled with the full-data experiment because they evaluate different training and test samples.
On full Superconductivity, the stage cap is selected in 3 of 5 Natural fits, 4 of 5 Ordinary fits and all 5 Homoscedastic fits, giving 12 of 15 overall.

\begin{table}[t]
\centering
\small
\setlength{\tabcolsep}{4pt}
\begin{tabular}{@{}lrrr@{}}
\toprule
Setting & Natural & Ordinary & Homosc. \\
\midrule
Airfoil, rows & 7.67 & 16.74 & 8.47 \\
Airfoil, conditions & 11.30 & 18.70 & 11.89 \\
Forest Fires & 5.35 & 5.38 & 5.36 \\
Parkinsons, subjects & 30.31 & 30.32 & 32.69 \\
Real Estate & 29.74 & 47.48 & 31.84 \\
Superconductivity, full & 34.55 & 115.79 & 43.47 \\
\bottomrule
\end{tabular}
\caption{Mean nominal 95\% prediction-interval widths, in each dataset's target units, including the transformed Forest Fires scale.}
\label{tab:new-widths}
\end{table}

\textbf{Synthetic dispersion and aggregate coverage.}
The 10 fresh synthetic replicates use independent base data and split draws, with regimes paired within each replicate.
Table~\ref{tab:synthetic-additional} reports sample standard deviations of replicate NLL and complements the main scale-recovery results with RMSE, coverage and width.
The mean oracle NLL references are 0.7477 in A, 0.9137 in B and 1.4408 in C.
They evaluate the true conditional mean and scale on the same finite test outcomes as each fitted model.
The earlier single-draw synthetic experiment is excluded from these summaries.

\begin{table}[t]
\centering
\footnotesize
\setlength{\tabcolsep}{3pt}
\begin{tabular}{@{}llrrrr@{}}
\toprule
Regime & Variant & NLL SD & RMSE & $C_{95}$ & Width \\
\midrule
A & Natural & 0.0348 & 0.6820 & 0.935 & 2.271 \\
& Ordinary & 0.0412 & 0.7029 & 0.924 & 2.302 \\
& Homosc. & 0.0324 & 0.6821 & 0.934 & 2.672 \\
\midrule
B & Natural & 0.0242 & 0.6181 & 0.943 & 2.392 \\
& Ordinary & 0.0294 & 0.6356 & 0.924 & 2.347 \\
& Homosc. & 0.0249 & 0.6167 & 0.948 & 2.391 \\
\midrule
C & Natural & 0.0339 & 1.3512 & 0.937 & 4.503 \\
& Ordinary & 0.0335 & 1.3936 & 0.927 & 4.592 \\
& Homosc. & 0.0318 & 1.3520 & 0.936 & 5.304 \\
\bottomrule
\end{tabular}
\caption{Additional synthetic results over 10 replicates per regime.
Except for NLL SD, entries are means of test metrics; coverage and width concern nominal 95\% prediction intervals.}
\label{tab:synthetic-additional}
\end{table}

\textbf{Paired synthetic comparisons.}
Table~\ref{tab:synthetic-paired-counts} counts strict NLL wins on the shared test outcomes within each replicate.
Natural beats Ordinary in all 10 replicates of every regime and Homoscedastic in every replicate of A and C.
Homoscedastic beats Natural in 9 of 10 B replicates.

\begin{table}[t]
\centering
\small
\setlength{\tabcolsep}{4pt}
\begin{tabular}{@{}lrrr@{}}
\toprule
Comparison & A & B & C \\
\midrule
Natural beats Ordinary & 10/10 & 10/10 & 10/10 \\
Natural beats Homosc. & 10/10 & 1/10 & 10/10 \\
Homosc. beats Natural & 0/10 & 9/10 & 0/10 \\
\bottomrule
\end{tabular}
\caption{Paired counts of lower test NLL across 10 replicates per regime.}
\label{tab:synthetic-paired-counts}
\end{table}

\textbf{Pointwise disagreements.}
Six explanations differ between the judge and annotator on at least one criterion.
Table~\ref{tab:judge-disagreement-counts} gives the direction of each criterion's disagreements.
Both evidence-support disagreements concern intended-faithful Real Estate descriptions.
Three jargon explanations account for all clarity disagreements and three of the four interval disagreements: the annotator rejects clarity but accepts interval consistency, while the judge does the reverse.
The annotator accepts clarity on two other jargon explanations, despite all five sharing the same final sentence.

\begin{table}[t]
\centering
\small
\setlength{\tabcolsep}{4pt}
\begin{tabular}{@{}lrr@{}}
\toprule
& Judge pass, & Judge fail, \\
Criterion & annotator fail & annotator pass \\
\midrule
Clarity & 3 & 0 \\
Interval consistency & 0 & 4 \\
Evidence support & 0 & 2 \\
\bottomrule
\end{tabular}
\caption{Directional disagreement counts out of 30 explanations per criterion.
These are comparisons with the observed annotation, not adjudicated error counts.}
\label{tab:judge-disagreement-counts}
\end{table}

\textbf{Pairwise order checks.}
Table~\ref{tab:judge-order-counts} distinguishes answer-identity consistency across swaps from literal position choices.
The five interval-misinterpretation edits all fail interval consistency individually but win the pairwise comparison when placed second.
These contrasts were checked against the raw responses using the frozen verdict parser, not inferred from the edit categories alone.

\begin{table}[t]
\centering
\small
\setlength{\tabcolsep}{4pt}
\begin{tabular}{@{}lr@{}}
\toprule
Check & Count \\
\midrule
Same underlying verdict after swap & 4/15 \\
Changed verdict after swap & 11/15 \\
First position preferred & 4/30 \\
Second position preferred & 20/30 \\
Explicit tie & 6/30 \\
Interval edit rejected individually & 5/5 \\
Interval edit preferred when second & 5/5 \\
\bottomrule
\end{tabular}
\caption{Pairwise consistency counts use 15 evidence cases; position counts use 30 displayed orders.
The final two rows concern the same five interval-error explanations.}
\label{tab:judge-order-counts}
\end{table}

\textbf{Annotation procedure.}
One annotator labelled all 30 explanations on the three criteria, using the rubric text supplied to the judge.
The annotation form contained an evidence packet, an explanation and fields for the three binary labels and a comment.
Items were presented individually, and the annotator's responses were recorded in a single sheet.
That sheet holds only an opaque item identifier, the evidence packet, the explanation, the three binary labels, an annotator identifier and a free-text note.
It carries no construction intent, no case identifier and no judge output.
Shuffling did not keep related items apart: 4 of the 15 explanation pairs are adjacent.
All 90 labels and all 30 notes are complete, and the evidence and explanation fields are unchanged from the archived blank sheet.
The completed sheet and its SHA-256 are preserved with the batch.

\textbf{Judge reproducibility.}
The judge weights are a third-party Q8\_0 quantisation of MiniCPM4.1-8B \cite{openbmb2025minicpm41}, identified by the Hugging Face repository, revision and file hash below.
\begingroup
\footnotesize
% A minipage keeps the hash from splitting across the column break.
\noindent\begin{minipage}{\columnwidth}
\begin{verbatim}
base    openbmb/MiniCPM4.1-8B
gguf    mradermacher/MiniCPM4.1-8B-GGUF
rev     7032ea8d4fc677dd3263c25d7c5343a8094c3151
sha256  9e939664e74b400b618d0e700a0794a2
        d58a889c487ec6c62983354b407bce7e
\end{verbatim}
\end{minipage}
\endgroup
\par\smallskip
\noindent Ollama~0.9.6 \cite{ollama2025} served sequential CPU requests with a 2{,}048-token context, 4 threads and at most 160 output tokens.
The ChatML prefix used non-reasoning mode, and JSON schemas constrained the verdict format; outputs are discrete labels rather than calibrated label probabilities.
The three-case development smoke test exposed order sensitivity; its prompt was retained unchanged for evaluation rather than retuned against evaluation results.
Raw prompts, responses, the frozen code and model configuration are preserved with checksums.
